\begin{abstract}

Deep learning models trained with mean squared error (MSE) on univariate continuous glucose monitoring (CGM) data produce predictions that are systematically delayed by approximately 50 to 55~minutes relative to the true signal. This time-shift artefact causes hypoglycaemia events, defined as glucose values falling below 70~mg/dL, to be detected too late or missed entirely. This thesis evaluates three strategies for mitigating the delay in a controlled ablation study on the T1DATA cohort of 36 adults with Type~1 diabetes, using LSTM and PatchTST architectures throughout.

Objective~1 replaces MSE with offset-aware loss functions: Bounded-Lag (tolerance $D{=}3$ steps) and Soft-DTW ($\gamma{=}1.0$). Objective~2 extends single-endpoint supervision to simultaneous prediction of all $H{=}12$ future steps. Objective~3 introduces event-centric evaluation via a tolerance-window sweep and trains dedicated binary classifiers for direct event detection. Exactly one component changes per objective; all other settings are fixed.

Multi-step MSE training achieves the best RMSE$^*$ for LSTM (20.74~mg/dL, versus 21.65~mg/dL for the MSE baseline). For PatchTST, the MSE baseline itself achieves the lowest RMSE$^*$ (14.21~mg/dL). Offset-aware loss functions increase RMSE$^*$ for both architectures despite reducing the mean prediction delay, indicating that temporal realignment is achieved at the cost of prediction shape accuracy. Loss function modifications alone reduce the mean prediction delay by at most 5.2~minutes without producing clinically meaningful improvements in event detection. Direct binary classifiers outperform all forecast-derived detectors under the recall-weighted F2 criterion: the LSTM classifier achieves F2~=~0.463 and recall~=~0.693, compared to a best forecast-derived F2 of 0.013, a large and statistically significant improvement after Holm--Bonferroni correction over a forecast-derived baseline that operates near the performance floor. A lead-time analysis qualifies this result: 49.5\% of the LSTM classifier's true-positive alarms are issued at zero lead time and therefore confirm an ongoing event rather than anticipate it. A performance inversion is observed across architectures: the LSTM, which is the weaker forecaster by RMSE$^*$, is the stronger classifier by F2. The PatchTST classifier achieves only F2~=~0.178, because its Reversible Instance Normalisation removes the absolute glucose level information that is essential for threshold-based classification.

These results indicate that glucose trajectory prediction and hypoglycaemia detection are best addressed as separate tasks. For a system requiring both capabilities, a multi-step MSE model for trajectory forecasting paired with a dedicated LSTM binary classifier is the strongest configuration evaluated here, although the lead-time analysis shows that its detection value is largely confirmatory rather than anticipatory and that clinical use would require further reduction of its false-alarm rate.

\end{abstract}
